%%
%% 2020 02 fp@bms-w.ch
%%

%% Quadratische Gleichungen: Substitution
\section{Substitution}\index{Substitution!quadratische Gleichungen}
\GESO{Kein Thema bei GESO ab RLP 2030!}

Zur Erinnerung beim Faktorisieren:

\TNT{2.8}{
$7x\cdot{}(\sqrt{x}+b)^2 + 5z\cdot{}(b+\sqrt{x})^2 = 7x\cdot{}u +
  5z\cdot{}u = (7x+5z)\cdot{}u = (7x+5z)\cdot{}(\sqrt{x}+b)^2$
}

Oft können kompliziertere Gleichungen mittels einer geeigneten \textbf{Substitution} (= Ersetzung\footnote{lat. \textbf{substitu\=o} = an die Stelle \textit{jmds. o. einer Sache} setzen.
})
in eine einfachere verwandelt werden. Wie \zB hier:

\newcommand\tmpPart{\left(\frac{1}{3}\cdot{}(5x+0.5)\right)}

$$\tmpPart^2 - 10 = 3\cdot{}{\tmpPart{}}$$

%%
\TRAINER{«Mit Farben geht alles besser.»\\}%%
\TNT{2.0}{$${\color{ForestGreen}\tmpPart^{\color{black}2}} -10 = 3{\color{ForestGreen}{\tmpPart{}}}$$}

\textbf{Substitution}

\TNT{1.2}{Wir ersetzen: $u := {\color{ForestGreen}\tmpPart{}}$}

und die ursprüngliche Gleichung ist nun äquivalent zu
\TNT{1.6}{$${\color{ForestGreen}u}^2 - 10 = 3{\color{ForestGreen}u}$$
  $$u^2 -3u - 10 = 0$$
}

\begin{bbwFillInTabular}{c|c|c}
  a&b&c\\
  \LoesungsRaum{1}&\LoesungsRaum{-3}&\LoesungsRaum{-10}
  \end{bbwFillInTabular}

\GESO{Substituierte Gleichung \GESO{mit TR }lösen:}
\TALS{Substituierte Gleichung mit $abc$-Formel:

\TNT{2}{
  $u_{1,2} = \frac{3\pm\sqrt{9+40}}{2} = \frac{3\pm 7}{2}$}%% END TNT
}%% END TALS


$$\LoesungsMenge{}_u=\LoesungsRaumLang{\{\LoesungsRaum{-2}; \LoesungsRaum{5}\}}$$

\TRAINER{\textbf{Nomenklatur:}\\
  Substitutionsbasis: $\tmpPart^2 - 10 = 3\tmpPart{}$\\
  Substituendum: $\tmpPart{}$\\
  Substituens: $u$
}%%

\newpage


\textbf{\TRAINER{Rücksubstitution}\noTRAINER{.............................................}}\index{Rücksubstitution}\\
Nach der ursprünglichen Variable ($x$) auflösen:


\TNTeop{
  $$u=\frac13 \cdot{} (5x+0.5) \Longleftrightarrow
  3u = 5x+0.5 \Longleftrightarrow
  3u-0.5 = 5x \Longleftrightarrow
  x = \frac{3u-0.5}{5}$$

Die gefundenen $u$-Werte (-2 und 5) einsetzen:

  
$$u_1 = -2: x_1 = \frac{3(-2)-0.5}{5} = \frac{-6.5}{5} = -1.3$$

$$u_2 =  5: x_2 = \frac{3( 5)-0.5}{5} = \frac{14.5}{5} =  2.9$$


$$\lx=\{\LoesungsRaum{-1.3}; \LoesungsRaum{2.9}\}$$}%% END TNTeop

%%%%%%%%%%%%%%%%%%%%%%%%%%%%%%%%%%%55555
\newpage

\subsection*{Aufgaben}

\olatLinkGESOKompendium{2.3.3}{17}{52}

%%\GESO{\AadBMTA{184}{29. a) 30. a)}}
\TALS{\AadBMTA{184}{29. a) c), 30. a) b)}}


  %% TALS hat hier ein einfacheres Beispiel, da TALS ohne TR lösen muss:
  \TALS{
\aufgabenFarbe{$$\left(3x-\frac15\right)^2-4\left(3x-\frac15\right)=21$$}%% end AufgabenFarbe
\TNTeop{
  Substitution:
  $$u:= 3x-\frac15$$
  $$u^2 -4u -21 = 0$$
  $$u_{1,2} = \frac{4\pm\sqrt{16+84}}{2}=\frac{4\pm 10}{2}\Rightarrow u_1 = 7; u_2 = -3$$

  Resub:

  $$ 7 = 3x-\frac15 \Rightarrow 7.2=3x \Rightarrow x = 2.4$$
  $$-3 = 3x-\frac15 \Rightarrow -2.8=3x \Rightarrow x =
  \frac{-28}{30} = \frac{-14}{15}$$
}
  }%% END TALS

  \newpage

\subsection{Ein weiteres Konzept: Potenzgleichungen}
  
  \aufgabenFarbe{Berufsmaturitätsprüfung 2020: Aufgabe 4 von Serie 2:
  $$(x-5)^4 - \frac{(x-5)^2}{3} = 8$$}%% END Aufgabenfarbe
  \TNT{18}{
    Substituiere: $u:=(x-5)^2$

    Substituierte Gleichung:

    $$u^2 -\frac{u}3 = 8$$

    In Grundform:

    $$u^2 -\frac13 \cdot{} u - 8 = 0$$

    Lösung (TR) $u$:   $$u_1=3  \text{ bzw. } u_2=-\frac83$$


  Rücksubstitution:

  $$u=(x-5)^2$$
  $u1:$ $$3=(x_1-5)^2$$
  Wurzel $\pm$ (!)ziehen.
  $$\pm\sqrt{3} = x_1-5$$
  
  $$5\pm\sqrt{3} = x_1$$
    $u_2$ ist Scheinlösung, denn $(x-5)^2 [=u]$ kann nicht negativ sein.
  \vspace{20mm}
  }%% END TNT
  
    $$\lx = \LoesungsRaumLang{\left\{5-\sqrt{3}; 5+\sqrt{3}\right\}}$$
  
  \newpage
  
  Die folgende Aufgabe stammt aus der \textit{Nullserie}-Maturaprüfung
  \TALS{Gesundheitlich-soziale Ausrichtung}:

  Lösen Sie mit Hilfe einer geeigneten Substitution \TRAINER{Tipp: $Z^{10} = \left(Z^5\right)^2$ }:
  $$\left(\frac{x}{3}-3.5 \right)^{10} -2\left( \frac{x}{3} - 3.5 \right)^5 =-1$$
  
  \TNTeop{Substituiere $u = \left(\frac{x}{3} - 3.5\right)^5$. Somit lautet die Gleichung
$$u^2 - 2u +1 = 0$$
    Was uns zu $u = 1$ bringt. Damit ist $1 = \left(\frac{x}{3} - 3.5 \right)^5$ und somit auch $1 = \frac{x}{3} - 3.5$. Auf beiden Seiten 3.5 addieren und danach mit 3 multiplizieren liefert $x = 13.5$.
 }%% End TNT
  
  %%%%%%%%%%%%%%%%%%%%%%%%%%%%%%%%%%%%

\subsection*{Aufgaben}
  

\DeclareRobustCommand{\tempA}{Aufgabe 14. a) b) c) d) h) i)}
\GESO{\olatLinkArbeitsblatt{GlQuad}{https://olat.bms-w.ch/auth/RepositoryEntry/6029794/CourseNode/111365559999287}{\tempA{}}}%% end GESO
\TALS{\olatLinkArbeitsblatt{GlQuad}{https://olat.bms-w.ch/auth/RepositoryEntry/6029786/CourseNode/111365560138736}{\tempA{}}}%% end TALS

\GESO{\olatLinkGESOKompendium{2.3.3.}{17}{53. und 54.}}

\TALS{\AadBMTA{184}{28. a), d)}}
\newpage

\subsection*{Referenzaufgabe Wurzelgleichung mit Substitution}
\TALS{
  \aufgabenFarbe{
    $$-6\sqrt{x-5} + (x-5)=-9$$
    }%% end Aufgabenfarbe
  \TNT{14}{
    $$u:=\sqrt{x-5}$$
    $$-6u+u^2=-9$$
    $$u^2-6u+9=0$$
    $$(u-3)^2=0$$
    $$\LoesungsMenge{}_u = \{3\}$$
    Resub:
    $$3=u=\sqrt{x-5}$$
    $$9=x-5$$
    $$\LoesungsMenge{}_x = \{14\}$$
  }%% END TNT
\newpage
}%% END TALS

\subsection*{Aufgaben}
\TALS{
  \aufgabenFarbe{
    $$\sqrt{x-1} + x-1 = 12$$
    }%% end Aufgabenfarbe
  \TNT{18}{

    $$u:= \sqrt{x-1}$$
    $$u + u^2 = 12$$
    $$u^2 + u - 12 = 0$$
    $$(u+4)(u-3)=0$$
    $$u_1=-4; u_2 = 3$$
    Rücksubst:

    $$x -1=u^2$$
    $$x = u^2 + 1$$
    $$u_1: x_1 = (-4)^2 + 1 = 17$$
    $$u_2: x_2 = (3)^2 + 1 = 10$$
    Probe: Nur $x_2 = 10$ funktioniert. $\lx = \{10\}$
  }%% END TNT
}%% END TALS



\TALS{\AadBMTA{197}{8. b, c)}}


\newpage


